% Pakiety LaTeX — kolejność ma znaczenie.
% Kompilacja: XeLaTeX + Biber (babel, polyglossia i fontspec wymagają silnika Unicode).

% --- Język ---------------------------------------------------------------
\usepackage[polish]{babel}
\usepackage{csquotes}

% --- Czcionki (fontspec wymaga XeLaTeX lub LuaLaTeX) ---------------------
\usepackage{fontspec}
\defaultfontfeatures{Ligatures=TeX}

% --- Geometria strony (odpowiednik for-printing z Typst) -----------------
\usepackage[
  a4paper,
  top=1in,
  bottom=1.25in,
  inner=2.25in,
  outer=1.25in,
  headheight=14pt
]{geometry}

% --- Interlinia i akapity ------------------------------------------------
\usepackage{setspace}
\onehalfspacing
\setlength{\parindent}{1em}
\setlength{\parskip}{0.2\baselineskip plus 0.1\baselineskip}

% --- Grafika, kolory -----------------------------------------------------
\usepackage{graphicx}
\usepackage[dvipsnames,table]{xcolor}
\graphicspath{{contents/assets/images/}}

% --- Tabele --------------------------------------------------------------
\usepackage{booktabs}
\usepackage{array}
\usepackage{longtable}

% --- Matematyka ----------------------------------------------------------
\usepackage{amsmath}
\usepackage{amssymb}

% --- Podpisy rysunków / tabel -------------------------------------------
\usepackage{caption}
\usepackage{subcaption}
\captionsetup{font=small,labelfont=bf}

% --- Listingi kodu (zamiennik codly) -------------------------------------
\usepackage{listings}
\definecolor{codebg}{HTML}{D2EBEB}
\definecolor{codeNumGray}{gray}{0.55}
\definecolor{codeKeyword}{HTML}{2E86AB}
\definecolor{codeComment}{HTML}{6A8759}
\definecolor{codeString}{HTML}{A9883E}

% --- Nagłówki i stopki ---------------------------------------------------
\usepackage{fancyhdr}
\usepackage{emptypage}

% --- Stylizacja rozdziałów / sekcji --------------------------------------
\usepackage{titlesec}

% --- Bibliografia --------------------------------------------------------
\usepackage[
  backend=biber,
  style=numeric,
  sorting=none,
  maxbibnames=99,
  giveninits=true,
  urldate=long
]{biblatex}
\addbibresource{bibliography/references.bib}

% --- Wsparcie dla pionowych odstępów -------------------------------------
\usepackage{needspace}

% --- Odnośniki (zawsze na końcu preambuły, przed cleveref jeśli używany) -
\usepackage[
  hidelinks,
  unicode,
  bookmarksnumbered,
  breaklinks
]{hyperref}

% --- Dodatkowe udogodnienia ---------------------------------------------
\usepackage{microtype}
\usepackage{enumitem}
\setlist{itemsep=0.2em,topsep=0.3em}
